\section{Related Work}
\label{sec:related works}
\subsection{Inference-Time Reward Alignment}
Sequential Monte Carlo (SMC)~\cite{doucet2001sequential, del2006sequential, chopin2020introduction} has proven effective in guiding the generation process of score-based generative models for inference-time reward alignment~\cite{cardoso2024mcgdiff, trippe2023smcdiff, dou2024fps, wu2023tds, kim2025DAS}. Prior SMC-based methods differ in their assumptions and applicability. For instance, FPS~\cite{dou2024fps} and MCGdiff~\cite{cardoso2024mcgdiff} are specifically designed for linear inverse problems and thus cannot generalize to arbitrary reward functions. SMC-Diff~\cite{trippe2023smcdiff} depends on the idealized assumption that the learned reverse process exactly matches the forward noising process—an assumption that rarely holds in practice. TDS~\cite{wu2023tds} and DAS~\cite{kim2025DAS} employ twisting and tempering strategies respectively to improve approximation accuracy while reducing the number of required particles. Despite these variations, all aforementioned SMC-based approaches share a common limitation that they initialize particles from the standard Gaussian prior, which is agnostic to the reward function. This mismatch can result in poor coverage of high-reward regions and reduced sampling efficiency.

In addition to multi-particle systems like SMC, single-particle approaches have also been explored for inference-time reward alignment~\cite{chung2023dps, yu2023freedom, bansal2024universal, ye2024tfg}. These methods guide generation by applying reward gradients along a single sampling trajectory. However, they are inherently limited in inference-time reward alignment, as simply increasing the number of denoising steps does not consistently lead to better sample quality. In contrast, SMC-based methods allow users to trade computational cost for improved reward alignment, making them more flexible and scalable in practice.

\subsection{Fine-Tuning-Based Reward Alignment}
Beyond inference-time methods, another line of work focuses on fine-tuning score-based generative models for reward alignment. Some approaches perform supervised fine-tuning by weighting generated samples according to their reward scores and updating the model to favor high-reward outputs~\cite{lee2023aligning, wu2023hps}, while others frame the denoising process as a Markov Decision Process (MDPs) and apply reinforcement learning techniques such as policy gradients~\cite{black2024ddpo} or entropy-regularized objectives to mitigate overoptimization~\cite{fan2023reinforcement, Wallace_2024_DPO, yang2-24D3po}. These RL-based methods are especially useful when the reward model is non-differentiable but may miss gradient signals when available. More recent methods enable direct backpropagation of reward gradients through the generative process~\cite{clark2024draft, prabhudesai2024alignprop}. Alternatively, several works~\cite{uehara2024finetuningcontinuoustimediffusionmodels, tang2024finetuningdiffusionmodelsstochastic, domingo-enrich2025adjoint} adopt a stochastic optimal control (SOC) perspective, deriving closed-form optimal drift and initial distributions using pathwise KL objectives. \\
While fine-tuning-based methods are an appealing approach, they have practical limitations in that they necessitate costly retraining whenever changes are made to the reward function or the pretrained model. Further, it has been shown that fine-tuning-based methods exhibit mode-seeking behavior~\cite{kim2025DAS}, which leads to low diversity in the generated samples.
